% Options for packages loaded elsewhere
\PassOptionsToPackage{unicode}{hyperref}
\PassOptionsToPackage{hyphens}{url}
\documentclass[
  ignorenonframetext,
]{beamer}
\newif\ifbibliography
\usepackage{pgfpages}
\setbeamertemplate{caption}[numbered]
\setbeamertemplate{caption label separator}{: }
\setbeamercolor{caption name}{fg=normal text.fg}
\beamertemplatenavigationsymbolsempty
% remove section numbering
\setbeamertemplate{part page}{
  \centering
  \begin{beamercolorbox}[sep=16pt,center]{part title}
    \usebeamerfont{part title}\insertpart\par
  \end{beamercolorbox}
}
\setbeamertemplate{section page}{
  \centering
  \begin{beamercolorbox}[sep=12pt,center]{section title}
    \usebeamerfont{section title}\insertsection\par
  \end{beamercolorbox}
}
\setbeamertemplate{subsection page}{
  \centering
  \begin{beamercolorbox}[sep=8pt,center]{subsection title}
    \usebeamerfont{subsection title}\insertsubsection\par
  \end{beamercolorbox}
}
% Prevent slide breaks in the middle of a paragraph
\widowpenalties 1 10000
\raggedbottom
\AtBeginPart{
  \frame{\partpage}
}
\AtBeginSection{
  \ifbibliography
  \else
    \frame{\sectionpage}
  \fi
}
\AtBeginSubsection{
  \frame{\subsectionpage}
}
\usepackage{iftex}
\ifPDFTeX
  \usepackage[T1]{fontenc}
  \usepackage[utf8]{inputenc}
  \usepackage{textcomp} % provide euro and other symbols
\else % if luatex or xetex
  \usepackage{unicode-math} % this also loads fontspec
  \defaultfontfeatures{Scale=MatchLowercase}
  \defaultfontfeatures[\rmfamily]{Ligatures=TeX,Scale=1}
\fi
\usepackage{lmodern}
\ifPDFTeX\else
  % xetex/luatex font selection
\fi
% Use upquote if available, for straight quotes in verbatim environments
\IfFileExists{upquote.sty}{\usepackage{upquote}}{}
\IfFileExists{microtype.sty}{% use microtype if available
  \usepackage[]{microtype}
  \UseMicrotypeSet[protrusion]{basicmath} % disable protrusion for tt fonts
}{}
\makeatletter
\@ifundefined{KOMAClassName}{% if non-KOMA class
  \IfFileExists{parskip.sty}{%
    \usepackage{parskip}
  }{% else
    \setlength{\parindent}{0pt}
    \setlength{\parskip}{6pt plus 2pt minus 1pt}}
}{% if KOMA class
  \KOMAoptions{parskip=half}}
\makeatother
\setlength{\emergencystretch}{3em} % prevent overfull lines
\providecommand{\tightlist}{%
  \setlength{\itemsep}{0pt}\setlength{\parskip}{0pt}}
\usepackage{bookmark}
\IfFileExists{xurl.sty}{\usepackage{xurl}}{} % add URL line breaks if available
\urlstyle{same}
\hypersetup{
  hidelinks,
  pdfcreator={LaTeX via pandoc}}

\author{\texorpdfstring{}{}}
\date{}

\begin{document}

\begin{frame}[fragile]{Abstract}
\protect\phantomsection\label{abstract}
This paper presents a convergence study of \textbf{fixed-step gradient
descent} on a convex quadratic, framed as the computational exemplar of
the \href{https://github.com/docxology/template}{Research Project
Template}. The implementation lives in
\texttt{projects/code\_project/src/optimizer.py}; experiments and
figures are orchestrated by
\texttt{projects/code\_project/scripts/optimization\_analysis.py} and
hydrated into the manuscript through
\texttt{scripts/z\_generate\_manuscript\_variables.py}, so tables and
prose track \texttt{output/data/optimization\_results.csv} after every
pipeline run.

We evaluate 6 step sizes from \(\alpha = 0.01\) to \(\alpha = 2.5\),
spanning conservative, near-optimal, aggressive, and divergent regimes
for a unit Hessian model. The build chain exercises template
infrastructure end-to-end: scientific helpers
(\texttt{infrastructure.scientific.stability},
\texttt{infrastructure.scientific.benchmarking}), validation, rendering
(\texttt{infrastructure/rendering/pdf\_renderer.py}), and reporting.
Accessibility-oriented plotting defaults (colourblind-safe palette,
300\,dpi exports) are centralized in \texttt{optimization\_analysis.py}.

Contributions are \textbf{methodological} and \textbf{architectural}. On
the methods side, we relate empirical iteration counts and error decay
to the scalar contraction factor \(\rho(\alpha) = |1-\alpha|\) and
document cases where runs hit \(N_{\max} = 1000\) before meeting the
gradient tolerance. On the architecture side, we demonstrate a zero-mock
test suite on project \texttt{src/} (see
\href{https://github.com/docxology/template/blob/main/projects/code_project/tests/test_optimizer.py}{test\_optimizer.py}),
automated six-figure analysis, and reproducibility metadata
(configuration hash, artifact counts) injected into Section 6.

\textbf{Results (this configuration):} 4 of 6 grid points report
\texttt{converged=True} in the CSV; non-convergent rows flag either slow
progress at small \(\alpha\) under the iteration cap or instability when
\(|1-\alpha| \geq 1\). The analytical minimizer remains \(x^\ast = 1.0\)
with \(f(x^\ast) = -0.5\) for the configured \((A,b)\).

\textbf{Keywords:} gradient descent, reproducible research, zero-mock
testing, scientific infrastructure, pipeline orchestration
\end{frame}

\end{document}
